\documentclass[conference,10pt]{IEEEtran}
\usepackage[T1]{fontenc}
\usepackage{newtxtext,newtxmath}
\usepackage{amsmath}
\usepackage{booktabs}
\usepackage{graphicx}
\usepackage{pgfplots}
\usepackage{url}
\usepackage{cite}
\pgfplotsset{compat=1.18}

\begin{document}

\title{Deep Q-Networks for 10$\times$10 Battleship:\\
Learning to Hunt Without the Rules of Placement}

\author{\IEEEauthorblockN{Le Cheng}
\IEEEauthorblockA{Khoury College of Computer Sciences\\
Northeastern University, Boston, MA, USA\\
cheng.l1@northeastern.edu}}

\maketitle

\begin{abstract}
Battleship is a partially observable game in which an agent must sink a hidden
fleet in as few shots as possible. Tabular Q-learning fails on it: on a 4$\times$4
board, a Q-table keyed on the visible board plays no better than random, and the
classic 10$\times$10 board has more than $3\times10^{10}$ ship layouts. We train
Deep Q-Networks (DQN) with a fully convolutional dueling network, Double DQN
targets and action masking, using 256 vectorized games to train each agent in
about 14 minutes on a laptop. On 4$\times$4, a DQN learned from the raw board
reaches 8.92 expected shots, statistically level with the probability-density
heuristic (8.86) and within 2\% of the exact optimum (8.75). On 10$\times$10,
evaluated on 2{,}000 held-out boards, the board-only DQN reaches 53.0 shots,
and adding the density map as an input plane reaches 45.8, 1.2 shots behind
probability density (44.6). Using the density as a fixed baseline under a learned
correction did not improve on it.
\end{abstract}

\section{Introduction}
In single-player Battleship, ships of lengths 5, 4, 3, 3 and 2 are placed
uniformly at random on a 10$\times$10 grid, and each shot reveals \emph{miss},
\emph{hit} or \emph{sunk}. The hidden layout makes this a POMDP~\cite{kaelbling1998}.
The strongest simple strategy is \emph{probability density}~\cite{berry2011}: count
every position of every remaining ship consistent with the board, and fire at the
cell covered most often. It is greedy and counts ships independently, so in
principle a learner that estimates the shots still needed could improve on it.

Our group's tabular agents show why function approximation is needed. A Q-table
keyed on the visible 4$\times$4 board visited 3.5 million states in a million
games and played at random level (13.9 shots against 14.1). Keying the table on
the exact belief (the set of consistent layouts) reached 8.97, but the belief
cannot be stored on 10$\times$10, which has 30{,}093{,}975{,}536 layouts.
Silver and Veness introduced 10$\times$10 Battleship as a POMDP planning
benchmark, solved online with Monte-Carlo tree search over particles~\cite{silver2010};
published learning approaches are hobby projects, the best DQN reaching about 52
shots~\cite{gaintel}.
We ask: \emph{can a DQN~\cite{mnih2015} learn 10$\times$10 Battleship without
being told the placement rules, and does the density map help it as an input
feature or as a starting policy?}

\section{Method}
\textbf{MDP.} The state $s$ is the visible board plus the lengths of the ships
still afloat. Observations are exact, so $s$ determines the belief and is a
Markov state. Actions are unfired cells. The reward is $-1$ per shot with
$\gamma=1$, so $Q(s,a)$ is minus the expected number of remaining shots and the
return equals minus the game length.

\textbf{Q-learning and DQN.} Q-learning~\cite{watkins1992} moves $Q(s,a)$
toward the target $y=r+\gamma\max_{a'}Q(s',a')$. DQN replaces the table with a
network $Q_\theta$ trained on the Huber loss of $Q_\theta(s,a)-y$, sampling
transitions from a replay buffer and computing $y$ with a periodically copied
target network $\theta^-$. We use Double DQN targets~\cite{vanhasselt2016}:
$y=-1+Q_{\theta^-}(s',\arg\max_{a'}Q_\theta(s',a'))$, and $y=-1$ when the game ends.
Fired cells are masked to $-\infty$ both when acting and inside the target.

\textbf{Input and network.} The board is encoded as planes of $10\times10$:
hit, miss, sunk, an all-ones on-board plane (so zero padding marks the edge),
and five planes counting the ships afloat of each length 1--5. The network has
four $3\times3$ convolutions with 32 channels and ReLU, giving a $9\times9$
receptive field, enough to see a carrier through any cell. A dueling
head~\cite{wang2016} computes $Q=V(s)+A(s,a)-\overline{A}(s)$: the advantage $A$ comes from a
$1\times1$ convolution (local), and the value $V$ from global mean and max
pooling (the shots still needed depend on the whole board). The network has
34.6k weights.

\textbf{Variants.} (A) \emph{board}: the planes above. (B) \emph{density}: plus
the probability-density map, normalized so the best cell is 1; it is a function
of the board, so it adds no information, only a feature. (C) \emph{residual}:
B with $Q=c\,\hat d(s,a)+Q_\theta(s,a)$, $c=1$, and $A$'s last layer initialized
to zero, so the untrained agent plays exactly like probability density and DQN
learns a correction~\cite{johannink2019}.

\textbf{Training.} 256 games run in parallel as tensors; one iteration fires one
shot in every game. Each transition is replayed 8 times on average (batch 256),
with a 200k-transition buffer, Adam (learning rate $3\times10^{-4}$), a target
copy every 1{,}000 steps, and $\varepsilon$ decaying linearly from 1.0 (C: 0.1)
to 0.02 over the first 10\% of 100k games. The games and the density map run
on the CPU and the network on the laptop GPU (Apple MPS). Unit tests check that
the vectorized games match the reference environment move for move, and that
the convolutional density equals the reference heuristic exactly.

\section{Experiments}
\textbf{Protocol.} Every 5k games we play the greedy policy on 500 validation
boards (seeds 100{,}000+) and keep the best checkpoint. For C the best checkpoint
was the untrained one, so we report its final checkpoint. The test set is 2{,}000
boards (seeds 0--1{,}999), the same for every agent, played through the
project's standard engine, which exposes only public information. We report means
with 95\% intervals and the \emph{paired} difference to probability density on
the same boards. On 4$\times$4 we play every one of the 264 layouts, giving the
exact expected shots of a deterministic policy.

\begin{table}[t]
\caption{Mean shots to sink the fleet (lower is better)}
\label{tab:results}
\centering
\setlength{\tabcolsep}{3.2pt}
\begin{tabular}{@{}lrrrr@{}}
\toprule
\multicolumn{5}{@{}l}{\emph{10$\times$10, 2{,}000 test boards}}\\
Agent & Shots & $\pm$95\% & Paired vs.\ density & Train \\
\midrule
Random & 95.27 & 0.22 & $+50.62\pm0.45$ & -- \\
Hunt-target + parity & 50.76 & 0.39 & $+6.11\pm0.55$ & -- \\
Probability density~\cite{berry2011} & 44.65 & 0.40 & -- & -- \\
Belief greedy (particles) & 44.50 & 0.41 & $-0.15\pm0.51$ & -- \\
DQN-A (board) & 53.00 & 0.45 & $+8.35\pm0.54$ & 12.4 min \\
DQN-B (+ density input) & \textbf{45.85} & 0.32 & $+1.20\pm0.46$ & 14.1 min \\
DQN-C (residual) & 46.05 & 0.32 & $+1.41\pm0.46$ & 14.0 min \\
\midrule
\multicolumn{5}{@{}l}{\emph{4$\times$4, all 264 layouts (exact for deterministic policies)}}\\
Random & 14.23 & & & \\
Q-table on visible board$^\dagger$ & 13.9 & & & \\
Q-table on exact belief & 8.97 & & & 7 min \\
DQN-A (board) & \textbf{8.92} & & $+0.06\pm0.26$ & 1.3 min \\
Probability density & 8.86 & & & \\
Optimal $V^*$ (exact) & 8.75 & & & \\
\bottomrule
\end{tabular}\\[2pt]
\footnotesize{$^\dagger$Measured by our group on 40 held-out layouts.}
\end{table}

\begin{figure}[t]
\centering
\begin{tikzpicture}
\begin{axis}[width=\columnwidth, height=4.3cm, xlabel={training games (thousands)},
  ylabel={validation shots}, ymin=42, ymax=64, xmin=0, xmax=100,
  legend columns=4, legend style={font=\scriptsize, at={(0.5,1.03)}, anchor=south, draw=none, fill=none,
  /tikz/every even column/.append style={column sep=4pt}},
  tick label style={font=\scriptsize}, label style={font=\small},
  legend cell align=left]
\addplot[thick, red!70!black, mark=*, mark size=0.9pt] table[x expr=\thisrow{games}/1000, y=eval_mean, col sep=comma] {data/curve_board.csv};
\addplot[thick, blue!70!black, mark=square*, mark size=0.9pt] table[x expr=\thisrow{games}/1000, y=eval_mean, col sep=comma] {data/curve_density.csv};
\addplot[thick, green!45!black, mark=triangle*, mark size=1pt] table[x expr=\thisrow{games}/1000, y=eval_mean, col sep=comma] {data/curve_residual.csv};
\addplot[dashed, gray, domain=0:100] {44.13};
\legend{A: board, B: + density input, C: residual, prob.\ density}
\end{axis}
\end{tikzpicture}
\caption{Greedy validation score during training on 10$\times$10 (500 boards;
untrained A and B start at 90.5 and 93.5, off the axis).}
\label{fig:curve}
\end{figure}

\textbf{Results.} Table~\ref{tab:results} and Fig.~\ref{fig:curve} summarize the
runs. On 4$\times$4, DQN-A learned from the raw board in 75 seconds and matched
probability density, while the board Q-table stayed at random level: the network's
generalization across boards is what the table lacked. On 10$\times$10, DQN-A
dropped from 90.5 to 52.4 validation shots within 15k games, comparable to the
$\sim$52 reported for an earlier DQN~\cite{gaintel}, then became unstable (53--62) while
its training loss rose steadily from 0.36 to 1.9. Adding the density
\emph{input} (B) gained 7.2 shots over A on the test set and came within
$1.20\pm0.46$ shots of probability density. The \emph{residual} agent (C) started
at the density level, but early updates pushed it to 56 shots; it recovered to
about 45 without ever beating its untrained starting point. No learned agent beat
probability density; belief greedy, the repo's strongest planner, is statistically
level with it ($-0.15\pm0.51$).

\textbf{Where shots are spent.} We split each game into shots fired while no
unsunk hit is showing (hunting) and while one is (targeting). Probability density
uses 27.5 and 17.2 shots. DQN-B uses only 17.1 for hunting but 28.8 while a
damaged ship is still afloat: it tends to leave a hit ship and probe elsewhere
before finishing it, which costs more than it saves.

\section{Discussion and Conclusion}
A convolutional DQN learns Battleship from experience at a fraction of the cost
of tabular methods: 4$\times$4 in about a minute and 10$\times$10 in about a
quarter of an hour, with no placement rules given. The density map is far more useful as
an input feature than as a fixed starting policy: as a feature the network can
use it selectively, while as a residual baseline early, poorly calibrated
corrections override a near-optimal policy. Greedy play is already within 2\% of
optimal on 4$\times$4, so the headroom above probability density on 10$\times$10
is likely below one shot, which a learner must capture without first losing
more.

\textbf{Limitations.} One training seed per variant; hyperparameters were not
tuned; checkpoints were selected on validation boards, disjoint from the test set.

\textbf{Future work.} A lower learning rate and less frequent target copies to
remove A's divergence, $n$-step returns and early stopping to need fewer games,
multiple seeds, and a residual agent that learns $V$ before it is allowed to
change the density's ranking.

\textbf{Reproducibility.} The code, unit tests, trained models and training
logs are in the project repository (package \texttt{battleship/dqn}); training
and benchmarking each take one command.

\section*{Acknowledgment}
The code, experiments and analysis were developed with the assistance of Claude
(Anthropic), an AI assistant. The author designed the study, directed the work
and is responsible for the results. The tabular Q-learning baselines are from the
author's project group.

\begin{thebibliography}{10}
\bibitem{mnih2015} V.~Mnih \emph{et al.}, ``Human-level control through deep
reinforcement learning,'' \emph{Nature}, vol.~518, no.~7540, pp.~529--533, 2015.
\bibitem{watkins1992} C.~J.~C.~H. Watkins and P.~Dayan, ``Q-learning,''
\emph{Machine Learning}, vol.~8, pp.~279--292, 1992.
\bibitem{vanhasselt2016} H.~van~Hasselt, A.~Guez, and D.~Silver, ``Deep
reinforcement learning with double Q-learning,'' in \emph{Proc.\ AAAI}, 2016.
\bibitem{wang2016} Z.~Wang \emph{et al.}, ``Dueling network architectures for
deep reinforcement learning,'' in \emph{Proc.\ ICML}, 2016.
\bibitem{kaelbling1998} L.~P. Kaelbling, M.~L. Littman, and A.~R. Cassandra,
``Planning and acting in partially observable stochastic domains,''
\emph{Artificial Intelligence}, vol.~101, pp.~99--134, 1998.
\bibitem{berry2011} N.~Berry, ``Battleship,'' DataGenetics, Dec.\ 2011.
[Online]. Available: \url{https://www.datagenetics.com/blog/december32011/}
\bibitem{silver2010} D.~Silver and J.~Veness, ``Monte-Carlo planning in large
POMDPs,'' in \emph{Advances in Neural Information Processing Systems 23}, 2010,
pp.~2164--2172.
\bibitem{gaintel} ``Deep reinforcement learning--of how to win at Battleship,''
GA-Intelligence. [Online]. Available:
\url{https://www.ga-intelligence.com/deep-reinforcement-learning-win-battleship}
\bibitem{johannink2019} T.~Johannink \emph{et al.}, ``Residual reinforcement
learning for robot control,'' in \emph{Proc.\ ICRA}, 2019.
\end{thebibliography}

\end{document}
